% Supplement subsections typed below as "Section~S<n>.<m>"; build.sh checks each number against the built supplement.
% supplement-ref supp:encoder=S2.3
% supplement-ref supp:learned-method=S6.1
\section{From Real-Valued Data to Permutations: Encoding and Multi-View Ensembles}\label{sec:encoding}

ArrowFlow ranks items, so real-valued features must first become rankings. Sorting a vector keeps only the order of its coordinates and throws away their sizes. For example, $[1.0,2.0,3.0]$ and $[0.01,100,100.01]$ give the same ranking. The encoder therefore first creates more coordinates to compare, by polynomial expansion and projection, and only then sorts. Several differently projected views then vote, and the experiments measure whether all this helps prediction.

\subsection{Encoding}\label{sec:encoder}

An encoder maps $x\in\R^d$ to a score vector $z\in\R^e$ and returns $\pi_x=\argsort(z)$, the ranking of the $e$ coordinate IDs by ascending score. It runs the four steps below, and every statistic they use is fitted on the training rows only.
\begin{enumerate}[leftmargin=2em,topsep=2pt,itemsep=1pt]
\item \emph{Imputation} replaces a missing value by the training mean of its column.
\item \emph{Polynomial expansion} of degree $p$ replaces $x$ by the vector $\phi(x)\in\R^{d'}$ of all monomials of degree at most $p$ in its coordinates, powers included. These monomials let the sorted scores depend on interactions between features, which changes the set of rankings that the encoder can produce from $\R^d$.
\item \emph{Standardization} centers every column and scales each nonconstant column to unit variance, giving $q(x)$.
\item \emph{Projection} multiplies by a matrix $W\in\R^{d'\times e}$, $z=q(x)W$, drawn by one of three strategies. A \emph{random} projection has independent standard normal entries. A \emph{target-aware} projection places coordinates of a linear discriminant analysis (LDA) before random coordinates. LDA is a linear map fitted on the training rows and their labels to separate the classes. A \emph{calibrated} projection is a random projection whose coordinates are standardized before sorting, so that no coordinate dominates through its variance.
\end{enumerate}
Section~S2.3 gives the exact rules of each step. At degree one the encoder also ignores scale. Sliding an example toward the training mean or away from it, along the ray from the mean through the example, leaves all of its rankings unchanged. Only its direction from the mean counts (Section~S2.3).

\subsection{Multi-View Ensemble}\label{sec:ensemble}

Different projection matrices rank the same input differently, and the ensemble uses this diversity (Figure~\ref{fig:pipeline}). A \emph{view} is one encoder together with the network trained on its rankings. ArrowFlow trains $K$ views independently, each from its own seed. View $k$ takes the projection strategy at position $k$ of the repeating cycle target-aware, random, calibrated.

Each view predicts with its nearest-neighbor readout (Section~\ref{sec:prediction}), and a \emph{majority vote} combines the views. The class predicted by the most views wins, even without an absolute majority, and the lowest class ID wins a tie. With the prototype readout, a \emph{Borda count} over the class rankings $\pi^{(L)}$ can replace the vote, as in classifier combination \citep{ho1994decision,vanerp2000variants}.

\subsection{Augmentation}\label{sec:augmentation}

Augmentation adds slightly perturbed copies of the training rankings, the permutation analog of small additive noise. Each copy swaps between $1$ and $s_{\max}$ random pairs of neighboring items, so it lies within footrule distance $2s_{\max}$ of its original. The originals and the validation rankings are kept unchanged. Whether augmentation is switched on is decided from the training rows alone (Section~S2.3).

\subsection{Settings}\label{sec:settings}

The encoded vocabulary size $e$, the degree $p$, the widths, the initial learning rate and the augmentation switch are set from the training rows alone. Some follow a fixed rule, and the others are chosen by an inner cross-validation.

Section~\ref{sec:protocol} describes the selection, and Section~S2 lists every setting with its symbol and value.

\subsection{An Encoder Trained by the Ranking Signal}\label{sec:learned-encoder}

The encoder above is fixed. Nothing in it changes when the network makes a mistake. This subsection replaces it by an \emph{encoder network}, a small neural network whose output scores are sorted exactly as before. It trains that network with motions of the same kind as those that train ArrowFlow's hidden layers. This is an extension of ArrowFlow, and it works in two stages. First, each view's encoder network is trained with its own layer of class filters. Second, the network is fixed and its class filters are set aside. ArrowFlow then trains its ranking layers on the rankings of the trained encoder, as it does with the fixed encoder. The signal that trains the encoder network also has an \emph{away term}, which pushes each coordinate away from the nearest wrong class filter. ArrowFlow's output vote never does that.

The difficulty is the sort. A neural network learns from gradients, but a ranking has no useful gradient. Nudging a score slightly either leaves the ranking unchanged or makes it jump. The motions offer another route (Figure~\ref{fig:target}). They say where each coordinate should sit in the ranking, and a position in the ranking corresponds to a score. So each motion can be translated into a target score, and the network can regress toward it.

The encoder network maps $x$ to scores $h\in\R^e$, and the input ranking is $\pi_x=\orderdesc(h)$, the $e$ coordinates in order of decreasing score. The fixed encoder of Section~\ref{sec:encoder} sorts in increasing order. Here the order is decreasing, so that position 1 holds the largest score. A ranking layer with one filter per class reads $\pi_x$, and its filters keep learning by the output layer's rule of Section~\ref{sec:training}. For an accepted example, the filter $r_y$ of the true class gives each coordinate a motion $m_y=m_{\pi_x\to r_y}$ by~\eqref{eq:motion}. The filter $r_w$ of the nearest wrong class, the other class whose filter lies nearest to $\pi_x$ by the footrule, gives it another motion $m_w=m_{\pi_x\to r_w}$. Coordinate by coordinate, $m_y$ is the motion that the class filter's vote returns (Section~\ref{sec:accumulation}), seen from the input's side and so with the opposite sign. Section~S6.1 gives the exact gating. The \emph{target conversion} turns these motions into numbers in two steps. First, it picks a target position. This is the coordinate's position in the true filter, minus half the displacement that the nearest wrong filter asks of it. Second, it reads the target score off the current scores. The target score is the score that now stands at the resulting position $q_p$,
\begin{equation}\label{eq:target}
\tilde h_{\pi_x[p]} = h_{\pi_x[q_p]},\qquad q_p = p+m_y[p]-\tfrac12\,m_w[p],
\end{equation}
with $q_p$ kept inside $[1,e]$ and the score interpolated between the two nearest positions. In position space, $q_p$ is one gradient step of unit size on a squared loss. It pulls each coordinate toward its position in the true filter and pushes it half as hard away from its position in the nearest wrong filter. Learning vector quantization moves its prototypes by a similar pull and push \citep{kohonen1990self}. The network then takes one step of Adam \citep{kingma2015adam}, a standard method of gradient-based training, on $\tfrac12\lVert h-\tilde h\rVert^2$, averaged over the accepted examples of the batch. The target is held fixed during the step. The sort itself is never differentiated. The discrete motion becomes a target in the network's own units, which is the idea of target propagation \citep{bengio2014auto,lee2015difference}.

% Figure: how the ranking signal trains the encoder network (Section 5.5). Drawn by hand; the scores in panel (b) are an
% illustration of Equation eq:target, not run data: e = 6, the coordinate at p = 5, its position 2 in the true filter
% (m_y = -3) and 6 in the nearest wrong filter (m_w = +1), so q_p = 5 - 3 - 1/2 = 1.5 and the target is (2.1 + 1.4)/2.
\begin{figure}[!htbp]
\centering
\resizebox{\textwidth}{!}{%
\begin{tikzpicture}[
    >=Stealth,
    box/.style={draw, rounded corners=2pt, minimum height=0.85cm, minimum width=1.7cm,
                font=\small, align=center, thick},
    databox/.style={box, fill=blue!8, minimum width=0.8cm},
    procbox/.style={box, fill=orange!12},
    sortbox/.style={box, fill=green!12},
    arr/.style={->, thick},
    lbl/.style={font=\footnotesize, text=black!70},
    slot/.style={draw, minimum width=1.05cm, minimum height=0.5cm, font=\small, inner sep=0pt},
    note/.style={font=\footnotesize, anchor=west, align=left},
]
% ------------------------------------------------ (a) the training loop
\node[font=\small\bfseries, anchor=west] at (-0.45,1.45) {(a) Training the encoder network};
\node[databox] (x) at (0,0) {$x$};
\node[procbox] (enc) at (2.4,0) {Encoder\\[-1pt]network};
\node[sortbox] (sort) at (6.0,0) {Sort,\\[-1pt]decreasing};
\node[sortbox] (filt) at (10.2,0) {Class filters\\[-1pt]$r_1,\dots,r_C$};
\draw[arr] (x) -- (enc);
\draw[arr] (enc) -- node[lbl, above=1pt] {scores $h$} (sort);
\draw[arr] (sort) -- node[lbl, above=1pt] {ranking $\pi_x$} (filt);
\node[lbl, font=\scriptsize\itshape] at (6.0,0.78) {never differentiated};
\node[procbox] (conv) at (8.6,-2.5) {Target conversion~\eqref{eq:target}};
\node[procbox] (step) at (2.4,-2.5) {Adam step on\\[-1pt]$\tfrac12\lVert h-\tilde h\rVert^2$};
\draw[arr] (filt.south) -- node[lbl, right] {motions $m_y$, $m_w$} ++(0,-1.0) -| ([xshift=7mm]conv.north);
\draw[arr, black!55] (sort.south) -- node[lbl, left] {current scores} ++(0,-1.0) -| ([xshift=-7mm]conv.north);
\draw[arr] (conv) -- node[lbl, above=1pt] {target scores $\tilde h$} node[lbl, below=1pt] {held fixed} (step);
\draw[arr, dashed] (step) -- node[lbl, left] {update} (enc);
% ------------------------------------------------ (b) one coordinate
\begin{scope}[xshift=13.9cm]
\node[font=\small\bfseries, anchor=west] at (-0.75,1.45) {(b) One coordinate, $e=6$};
\node[lbl] at (-0.35,0.62) {position};
\node[lbl] at (0.9,0.62) {score};
\foreach \p/\s in {1/2.1, 2/1.4, 3/0.9, 4/0.3, 5/{-0.4}, 6/{-1.2}} {
  \node[font=\small] at (-0.35,0.75-0.75*\p) {\p};
  \node[slot] (s\p) at (0.9,0.75-0.75*\p) {$\s$};
}
\node[slot, fill=red!15] at (s5) {$-0.4$};
% to the position in the true filter, then half a step away from the wrong filter
\draw[arr, blue!70!black] (s5.east) to[out=0, in=0, looseness=1.4] (s2.east);
\draw[red!75!black, line width=1.6pt] (0.2,-0.375) -- (1.6,-0.375);
\node[note, text=red!75!black] at (2.75,-0.375) {target position $q_p=2-\tfrac12=1.5$\\target score $\tilde h=1.75$};
\node[note, text=blue!70!black] at (2.75,-1.9) {true filter: position 2,\\so $m_y[p]=-3$};
\node[note] at (2.75,-3.0) {the coordinate, now at $p=5$};
\node[note, text=orange!60!black] at (2.75,-3.95) {wrong filter: position 6,\\so $m_w[p]=+1$};
\end{scope}
\end{tikzpicture}%
}
\caption{\textbf{How the ranking signal trains the encoder network.} (a)~The encoder network turns an example $x$ into scores $h$, and a sort in decreasing order gives the input ranking $\pi_x$. A ranking layer with one filter per class reads $\pi_x$ and keeps learning by its own vote. For an accepted example, the filter of the true class and the filter of the wrong class nearest to $\pi_x$ give the motions $m_y$ and $m_w$. The target conversion~\eqref{eq:target} turns them into target scores $\tilde h$. With $\tilde h$ held fixed, an Adam step on $\tfrac12\lVert h-\tilde h\rVert^2$ updates the network (dashed arrow). No gradient passes through the sort. After training, each view's encoder network is fixed and replaces that view's fixed encoder. (b)~The conversion for one coordinate, with $e=6$ scores listed from the largest down. The coordinate sits at position $p=5$ with score $-0.4$. The true filter puts it at position 2, so $m_y[p]=-3$ (blue arrow), and the nearest wrong filter puts it at position 6, so $m_w[p]=+1$. Its target position is $q_p=5-3-\tfrac12=1.5$ (red line). Its target score, $1.75$, lies halfway between the scores at positions 1 and 2, and the loss pulls its score toward that value. The scores are made up for the example.}
\label{fig:target}
\end{figure}


The encoder network has three fully connected layers of width $e$, and the first two start with weights drawn uniformly from 0 to 1. It trains for 60 passes over the training rows, with no validation checkpoint. Section~\ref{sec:learned} evaluates the extension, and Section~S6.1 gives every setting.
